% SC1004 Linear Algebra. Print on one A4 sheet, landscape, double-sided.
\documentclass[10pt,a4paper,landscape]{article}

\usepackage[T1]{fontenc}
\usepackage{lmodern}
\usepackage[margin=7.5mm]{geometry}
\usepackage{amsmath,amssymb,array,xcolor}

\renewcommand{\familydefault}{\sfdefault}
\definecolor{ink}{HTML}{19252C}
\definecolor{accent}{HTML}{173E50}
\definecolor{rulegray}{HTML}{A7B7BF}
\definecolor{wash}{HTML}{EDF2F4}
\pagestyle{empty}
\setlength{\parindent}{0pt}
\setlength{\parskip}{1.1pt}
\setlength{\fboxsep}{2.8pt}
\setlength{\tabcolsep}{3pt}
\setlength{\emergencystretch}{1em}
\raggedbottom

\newlength{\sheetcol}
\setlength{\sheetcol}{\dimexpr(\textwidth-6mm)/3\relax}
\newcommand{\Col}{\operatorname{Col}}
\newcommand{\Row}{\operatorname{Row}}
\newcommand{\Null}{\operatorname{N}}
\newcommand{\RREF}{\operatorname{rref}}
\newcommand{\rank}{\operatorname{rank}}
\newcommand{\spn}{\operatorname{span}}
\newcommand{\proj}{\operatorname{proj}}
\newcommand{\tr}{\operatorname{tr}}
\newcommand{\diag}{\operatorname{diag}}
\newcommand{\RR}{\mathbb{R}}
\newcommand{\norm}[1]{\lVert #1\rVert}
\newcommand{\abs}[1]{\lvert #1\rvert}
\newcommand{\smat}[1]{\bigl[\begin{smallmatrix}#1\end{smallmatrix}\bigr]}
\newcommand{\formula}[1]{\par\noindent\(\displaystyle #1\)\par}
\newcommand{\topic}[1]{%
  \par\addvspace{2.8pt}{\color{accent}\bfseries #1\par}\nobreak\vspace{-1pt}}
\newcommand{\checkstrip}[1]{%
  \par\vspace{2.5pt}%
  \noindent\colorbox{wash}{\parbox{\dimexpr\textwidth-2\fboxsep\relax}{\cue{True/false checks}\quad #1}}\par}
\newenvironment{panel}[1]{%
  \begin{minipage}[t]{\sheetcol}\vspace{0pt}%
  \raggedright\setlength{\parskip}{1.1pt}%
  {\color{accent}\bfseries\fontsize{10.5}{12}\selectfont \strut #1\par}%
  \vspace{.8pt}{\color{rulegray}\hrule height .45pt}\vspace{1pt}%
}{\end{minipage}}
\newcommand{\cue}[1]{\textbf{#1}}
\newcommand{\rowend}{\tabularnewline[1pt]}

% Use the standard fonts on full TeX installations. A partial macOS
% XeTeX cache can use local STIX glyphs for the missing math fonts.
\newif\ifpartialfontcache
\ifdefined\XeTeXversion
  \IfFileExists{lmex10.pfb}{}{%
    \IfFileExists{/System/Library/Fonts/Supplemental/STIXTwoMath.otf}{%
      \partialfontcachetrue
    }{}%
  }
\fi
\ifpartialfontcache
\makeatletter
\input{ot1lmr.fd}
\input{omllmm.fd}
\input{omslmsy.fd}
\makeatother
\DeclareFontShape{OT1}{lmr}{m}{n}{<->rm-lmr9}{}
\DeclareFontShape{OML}{lmm}{m}{it}{<-6.5>lmmi5 <6.5-10>lmmi9 <10->lmmi10}{}
\DeclareFontShape{OMS}{lmsy}{m}{n}{<-6.5>lmsy5 <6.5-10>lmsy8 <10->lmsy10}{}
\font\sheetSTIX="STIX Two Math" at 8pt
\font\sheetSTIXDisplay="STIX Two Math" at 10.5pt
\font\sheetSTIXBrackets="STIX Two Math" at 15pt
\renewcommand{\RR}{\mathord{\hbox{\sheetSTIX\char"211D}}}
\renewcommand{\sum}{\mathop{\mathchoice{\vcenter{\hbox{\sheetSTIXDisplay\char"2211}}}{\vcenter{\hbox{\sheetSTIX\char"2211}}}{\hbox{\sheetSTIX\char"2211}}{\hbox{\sheetSTIX\char"2211}}}\displaylimits}
\renewcommand{\prod}{\mathop{\mathchoice{\vcenter{\hbox{\sheetSTIXDisplay\char"220F}}}{\vcenter{\hbox{\sheetSTIX\char"220F}}}{\hbox{\sheetSTIX\char"220F}}{\hbox{\sheetSTIX\char"220F}}}\displaylimits}
\renewcommand{\smat}[1]{\mathopen{\vcenter{\hbox{\sheetSTIXBrackets[}}}\begin{smallmatrix}#1\end{smallmatrix}\mathclose{\vcenter{\hbox{\sheetSTIXBrackets]}}}}
\fi

% SC1004, one double-sided A4 sheet. Course guides and Finals Quick Reference v10 set the scope.
% Studocu theorem cross-check: https://www.studocu.com/sg/document/nanyang-technological-university/linear-algebra-for-computing/sc1004-final-exam-cheatsheet-linear-algebra-concepts-and-theorems/146017236
% Cross-checks: https://www.cs.cornell.edu/~tomf/notes/cs421-cheat-sheet.pdf
% https://ocw.mit.edu/courses/18-06sc-linear-algebra-fall-2011/5500c16aff18ebcc312df763e8f85992_MIT18_06SCF11_Ses2.12sum.pdf
% https://see.stanford.edu/materials/lsoftaee261/Formulas.pdf
\begin{document}
\fontsize{8.85}{10.1}\selectfont
\color{ink}
\noindent\begin{panel}{Systems, rank and bases}
\topic{Consistency and solution-set theorems}
Let $A=[a_1\,\cdots\,a_n]$ be $m\times n$ and $r=\rank A$.
$Ax=\sum_j x_j a_j$, so solving for $x$ tests membership in
the span of the columns. Use \cue{rref} on $[A\mid b]$.

\cue{REF and symbolic pivots}
Zero rows go last. Each leading entry is strictly to the right
of the one above. In RREF, each pivot is $1$ and alone in its column.
Before dividing by a parameter expression, split into its zero
and nonzero cases, then count pivots in each case.

\begin{tabular}{@{}>{\raggedright\arraybackslash}p{.40\linewidth}>{\raggedright\arraybackslash}p{.57\linewidth}@{}}
\cue{What you find} & \cue{What it tells you}\rowend
Pivot in the $b$ column & Inconsistent. No exact solution.\rowend
Consistent, $r=n$ & One solution.\rowend
Consistent, $r<n$ & Infinitely many. $n-r$ free variables.\rowend
Pivot in every column of $A$ & Columns LI. $x\mapsto Ax$ is one-to-one.\rowend
Pivot in every row of $A$ & Columns span $\RR^m$. Map is onto.\rowend
\end{tabular}

\formula{Ax=b\text{ is consistent}\iff b\in\Col A\iff\rank[A\mid b]=r.}
Given one solution $x_p$, all solutions are $x_p+\Null(A)$.
For $b=0$, the solution set is the null space (kernel).

\topic{Rank--nullity and the four fundamental spaces}
$\rank A+\dim\Null(A)=n$ for $n$ columns.
\begin{tabular}{@{}p{.27\linewidth}p{.10\linewidth}>{\raggedright\arraybackslash}p{.57\linewidth}@{}}
\cue{Space} & \cue{Dim.} & \cue{Build a basis}\rowend
$\Col A\subseteq\RR^m$ & $r$ & Pivot columns from the \cue{original} $A$.\rowend
$\Row A\subseteq\RR^n$ & $r$ & Nonzero echelon rows.\rowend
$\Null(A)\subseteq\RR^n$ & $n-r$ & Solve $Ax=0$. Give each free variable a turn at $1$ with the others at $0$.\rowend
$\Null(A^\top)\subseteq\RR^m$ & $m-r$ & Same null-space recipe on $A^\top$.\rowend
\end{tabular}

$\rank A=\rank A^\top\le\min(m,n)$. Row reduction preserves
row and null spaces. If \cue{rref} rejects a tall matrix,
append zero columns until it is square, then read the original block.

\topic{Subspace and spanning set theorems}
Test $0\in W$ and $\alpha u+\beta v\in W$ for all $u,v\in W$
and scalars. Any span is a subspace. Test span membership
with $[v_i]c=b$. LI means $[v_i]c=0$ forces $c=0$.
A basis is LI and spans. Drop a redundant vector without
changing the span. Keeping $u$, replacing $v$ by $v+cu$
also preserves the span. In dimension $d$, more than $d$
vectors or a set containing $0$ is dependent. Fewer than $d$
cannot span. With $d$ vectors, either basis check suffices.
For polynomials of degree $\le d$, $\dim P_d=d+1$. $\dim M_{m\times n}=mn$ and
$\dim\{S=S^\top\}=\dfrac{n(n+1)}{2}$ for $n\times n$ matrices.

\topic{Use structure to get rank quickly}
$\rank(AB)\le\min(\rank A,\rank B)$.
Multiplying by an invertible matrix preserves rank.
For $uv^\top$ with $u,v\ne0$, the rank is $1$, the columns
span $u$, and the rows span $v^\top$.
\end{panel}%
\hfill
\begin{panel}{Invertibility, determinants and LU}
\topic{Invertible Matrix Theorem (square \(A\))}
\formula{\begin{aligned}
 A\text{ invertible}&\iff\det A\ne0\iff\rank A=n\\
 &\iff\RREF(A)=I\iff\Null(A)=\{0\}\\
 &\iff\text{columns (or rows) form a basis}\\
 &\iff 0\text{ is not an eigenvalue}\\
 &\iff Ax=b\text{ has a unique solution for every }b.
\end{aligned}}
For square matrices, even a one-sided inverse suffices.
If $\det A=0$, any attainable $b$ has infinitely many preimages.
Use $\rank[A\mid b]$ to check the particular $b$.

\topic{Determinant laws and volume scaling}
Use \cue{det} for numeric determinants.
Swapping two rows or columns flips the sign. Scaling one
by $c$ multiplies the determinant by $c$. Adding a multiple
of another preserves it. A triangular determinant is the
product of the diagonal entries.

\formula{\begin{aligned}
 \det(AB)&=\det A\,\det B,\quad \det A^\top=\det A\\
 \det(cA)&=c^n\det A,\quad \det(A^k)=(\det A)^k\\
 \det(A^{-1})&=\frac{1}{\det A} \quad(A\text{ invertible}).
\end{aligned}}
For a rank-one change of the identity,
$\det(I+uv^\top)=1+v^\top u$.
The area or volume scale is $\abs{\det A}$.
Put edge vectors in the columns of $A$ to get the area
of a parallelogram or volume of a parallelepiped.
Stretching the unit disk by $\diag(a,b)$ gives an ellipse
of area $\pi ab$ for $a,b>0$. A negative determinant reverses orientation.

\topic{Read matrix products}
$(m\times n)(n\times p)$ gives $m\times p$.
Column $j$ of $AB$ is $Ab_j$, where $b_j$ is column $j$ of $B$.
Thus $AB=0\iff\Col B\subseteq\ker A$.
Use matrix keys for numerical products and inverses.
\formula{(AB)^\top=B^\top A^\top,\qquad
  (AB)^{-1}=B^{-1}A^{-1}.}
For compatible products, $\tr(AB)=\tr(BA)$.
Use $\tr(A+B)=\tr A+\tr B$ for square matrices of the same size. Also $(A^\top)^{-1}=(A^{-1})^\top$ when invertible.

\topic{Elementary matrices and LU factorisation}
One row operation on $I$ gives $E$. Then $EA$ performs the
same operation on $A$. Its inverse undoes the operation.

\cue{LU setup} $PA=LU$, where $P$ records row swaps, $L$ has
unit diagonal and stores elimination multipliers, and $U$
is the result of elimination. With no swaps $P=I$.
\formula{Ly=Pb,\qquad Ux=y.}
Use forward substitution, then back substitution. Reuse
the factors for a new $b$. For square $A$ with $s$ swaps,
$\det A=(-1)^s\prod_i u_{ii}$.

\cue{Conditional cancellation}
If $A$ is invertible, $\Null(AB)=\Null(B)$.
If $C$ is invertible, $\Null(B)=C\,\Null(BC)$.
\end{panel}%
\hfill
\begin{panel}{Geometry and linear maps}
\topic{Point--normal theorem}
Line $x=p+td$. Plane $x=p+su+tv$ with independent
directions. With normal $a\ne0$, use $a^\top(x-p)=0$.
A whole line $q+td$ lies in that plane iff
$a^\top(q-p)=0$ and $a^\top d=0$.
For a plane from two directions in $\RR^3$, take $a=u\times v$.
For two intersecting planes, use their normals' cross product
as a line direction.
\formula{u\times v=(u_2v_3-u_3v_2,\ u_3v_1-u_1v_3,\ u_1v_2-u_2v_1).}
$\norm{u\times v}=\norm u\norm v\sin\theta$.
A unit vector along a nonzero $d$ is $\dfrac{d}{\norm d}$.

\topic{Cauchy--Schwarz and triangle inequalities}
$\norm u^2=u^\top u$ and $d(u,v)=\norm{u-v}$.
For nonzero vectors,
\formula{\cos\theta=\frac{u^\top v}{\norm u\,\norm v}.}
Zero dot product means orthogonal.
$\abs{u^\top v}\le\norm u\norm v$ with equality iff $u,v$
are dependent. Also $\norm{u+v}\le\norm u+\norm v$.
$\norm{cu}=\abs c\,\norm u$.
For $u\perp v$, $\norm{u+v}^2=\norm u^2+\norm v^2$. In general, $\norm{u\pm v}^2=\norm u^2+\norm v^2\pm2u^\top v$.
\formula{d(y,\text{plane through }p)=\frac{\abs{a^\top(y-p)}}{\norm a}.}

\topic{Standard Matrix Theorem}
Test $T(\alpha u+\beta v)=\alpha T(u)+\beta T(v)$.
For linear $T$, $A=[T(e_1)\,\cdots\,T(e_n)]$ gives $T(x)=Ax$.
If $T(v_i)=w_i$, form $V=[v_i],W=[w_i]$.
Then $AV=W$, and $A=WV^{-1}$ when $V$ is a basis matrix.
In $AB$, apply $B$ first.

\topic{Kernel and image}
For linear $T:V\to W$, the kernel and image are subspaces
of the domain and codomain, respectively.
\formula{\ker T=\{v\in V:T(v)=0\},\quad \operatorname{im}T=\{T(v):v\in V\}.}
$\dim\ker T+\dim\operatorname{im}T=\dim V$ in finite dimension.
$T(u)=T(v)\iff u-v\in\ker T$. Hence one-to-one iff $\ker T=\{0\}$,
and onto iff $\operatorname{im}T=W$.
For $T(x)=Ax$, solve $Ax=0$ for $\ker T=\Null(A)$.
Then $\dim\ker T=n-r$ and $\operatorname{im}T=\Col A$ has dimension $r$.

\cue{Common maps}
Anticlockwise planar rotation is
$\smat{\cos\theta&-\sin\theta\\\sin\theta&\cos\theta}$.
Uniform scaling is $cI$. Horizontal shear is $\smat{1&k\\0&1}$. If $P$ projects orthogonally onto
a subspace, reflection across it is $2P-I$.

\topic{Change of Basis Theorem}
Let $P_{\mathcal B}=[b_1\,\cdots\,b_n]$ for a basis of $\RR^n$.
$x=P_{\mathcal B}[x]_{\mathcal B}$.
Solve for the weights to get $[x]_{\mathcal B}$.
For polynomials use coefficient columns in $1,t,t^2,\dots$.
For matrices list entries in a fixed order. Apply the same tests.
\formula{[x]_{\mathcal C}=P_{\mathcal C}^{-1}P_{\mathcal B}[x]_{\mathcal B}.}
Column $i$ of $P_{\mathcal C}^{-1}P_{\mathcal B}$ is $[b_i]_{\mathcal C}$.
\formula{[T]_{\mathcal C\leftarrow\mathcal B}=P_{\mathcal C}^{-1}AP_{\mathcal B}.}
Read right to left. Convert input, apply $A$, then convert output.
Using the same basis at both ends gives a similar matrix.
\end{panel}

\checkstrip{%
Row-equivalent matrices can have different column spaces.\quad
$\det(A+B)$ can differ from $\det A+\det B$.\quad
$T(0)=0$ alone leaves linearity undecided.}

\newpage
\setlength{\sheetcol}{\dimexpr(\textwidth-6mm)/2\relax}
\fontsize{10.5}{12.3}\selectfont
\noindent\begin{panel}{Projection, QR and least squares}
\topic{Orthogonal Decomposition Theorem}
$\Col(A)^\perp=\Null(A^\top)$, $\Row(A)^\perp=\Null(A)$,
and $\dim W^\perp=n-\dim W$ in $\RR^n$.
For an orthogonal basis $(u_i)$ of $W$,
\formula{c_i=\frac{u_i^\top y}{u_i^\top u_i},\quad
  \hat y=\proj_W y=\textstyle\sum_i c_i u_i.}
The signed component along $u$ is $\dfrac{u^\top y}{\norm u}$.
The split $y=\hat y+e$ with $\hat y\in W$, $e\in W^\perp$
is unique. For every $w\in W$,
\formula{\norm{y-w}^2=\norm e^2+\norm{\hat y-w}^2.}
\cue{Best Approximation Theorem} $\hat y$ is the unique
closest point in $W$ and $\norm e$ is the distance.
For $p+W$, project $y-p$ and add $p$.

\topic{Gram--Schmidt and the QR Factorisation Theorem}
Nonzero orthogonal vectors are LI. $Q^\top Q=I$ tests
orthonormal columns. $Q$ preserves dot products. When square,
$Q^{-1}=Q^\top$ and $\abs{\det Q}=1$.
\formula{v_j=a_j-\textstyle\sum_{i<j} q_i(q_i^\top a_j),\quad
  q_j=\frac{v_j}{\norm{v_j}}.}
Keep nonzero remainders to span the same space.
For $n$ LI columns, reduced $A=QR$ has $Q$ of size
$m\times n$ and upper triangular $R=Q^\top A$ of size $n\times n$.
$r_{ij}=q_i^\top a_j$ and $r_{jj}=\norm{v_j}>0$.
Full QR uses square $Q$ and appends zero rows to $R$ here.
\cue{Dependent columns} QR still exists. A square $R$ in
reduced QR is singular, so $Rx=c$ cannot have a unique solution.
$QQ^\top$ projects onto $\Col Q$. It projects onto $\Col A$
exactly when $\Col Q=\Col A$. Choose an orthonormal basis of $\Col A$ if needed.

\topic{Build and recognise a projector}
$P^2=P$ fixes its range. Add $P^\top=P$ to characterise
an orthogonal projector. For an orthonormal basis of $W$,
$P_W=QQ^\top$. For an LI spanning matrix $A$,
$P_W=A(A^\top A)^{-1}A^\top$.
For a line spanned by $u\ne0$, use $P_W=\dfrac{uu^\top}{u^\top u}$.
$P_Wy=\hat y$. The complementary projector $I-P_W$ maps
onto $W^\perp$ and gives $e=(I-P_W)y$.

\topic{Least Squares Theorem and normal equations}
A least-squares minimiser exists for every $b$.
It minimises $\norm{b-Ax}^2$ exactly when its residual
$e=b-A\hat x$ is in $\Null(A^\top)=(\Col A)^\perp$.
\formula{A^\top e=0\iff A^\top A\hat x=A^\top b.}
Solve the normal equations by elimination. With LI columns,
$\hat x=(A^\top A)^{-1}A^\top b$ or use reduced QR and solve
$R\hat x=Q^\top b$. All minimisers are $\hat x+\Null(A)$.
The fit $\hat b=A\hat x$ is always unique. Coefficients are
unique iff the columns are LI. Exact fit iff $e=0$.
Error is $\norm e$ and squared error is $\norm e^2$.
For $y\approx\sum_j c_j\phi_j(t)$, build $A_{ij}=\phi_j(t_i)$,
$b_i=y_i$, and solve the normal equations for $c$.

\cue{Gram matrix} $G=A^\top A$ is positive semidefinite since
$x^\top Gx=\norm{Ax}^2$. $\Null(G)=\Null(A)$, so $G$ is
positive definite and invertible iff the LS coefficients are unique.
\end{panel}%
\hfill
\begin{panel}{Real eigenvalues, diagonalisation and powers}
\topic{Eigenspaces are kernels}
Use \cue{EIGEN} for numerical eigenpairs.
Test a vector by $Av=\lambda v$, $v\ne0$. Test a value by
$\det(A-\lambda I)=0$. Eigenspace $E_\lambda=\ker(A-\lambda I)$.
Triangular $A$ has its eigenvalues on the diagonal.
For symbolic $2\times2$ work,
\formula{\det(tI-A)=t^2-(\tr A)t+\det A.}
\cue{AM} counts copies of $\lambda$ as a root.
\cue{GM} is $\dim E_\lambda=n-\rank(A-\lambda I)$.
For each eigenvalue, $1\le\mathrm{GM}\le\mathrm{AM}$.
Use the shifted rank test when roots repeat.
\cue{Similarity} preserves the characteristic polynomial,
rank and each eigenspace dimension.
If $B=S^{-1}AS$ and $v$ is an eigenvector of $A$ for $\lambda$,
then $S^{-1}v$ is an eigenvector of $B$ for $\lambda$.
When all roots are real, $\tr A=\sum_i\lambda_i$ and
$\det A=\prod_i\lambda_i$, counting algebraic multiplicity.

\topic{Diagonalisation Theorem}
$A$ is diagonalizable over $\RR$ iff it has $n$ LI real eigenvectors.
If there are $n$ real characteristic roots counting multiplicity,
test $\mathrm{AM}=\mathrm{GM}$ at each root. Eigenvectors for distinct
eigenvalues are LI, so $n$ distinct real eigenvalues suffice. Put eigenvectors in $P$ and
the matching eigenvalues in $D$.
\formula{AP=PD,\quad A=PDP^{-1},\quad A^k=PD^kP^{-1}.}
For $x_{k+1}=Ax_k$, set $c=P^{-1}x_0$ and use
$x_k=\sum_i c_i\lambda_i^k v_i$.
Active modes with $\abs\lambda<1$ vanish, $\lambda=1$ persists,
$\lambda=-1$ alternates, and $\abs\lambda>1$ grows in magnitude.
If all characteristic roots are real and have $\abs\lambda<1$, then $A^k\to0$.
If diagonalizable with only $1$ on the unit circle and all
others inside, retain the component for $\lambda=1$.

\cue{Reuse an eigenpair} $p(A)v=p(\lambda)v$.
For $A^{-1}$ the eigenvalue is $\dfrac{1}{\lambda}$ if invertible.
If $A^p=0$, all eigenvalues vanish and
$(I-A)^{-1}=I+A+\cdots+A^{p-1}$.

\topic{Infer from matrix structure}
\cue{Spectral Theorem} Real $S$ is symmetric iff $S=QDQ^\top$ for orthogonal $Q$ and real diagonal $D$.
For symmetric $S$, positive definite means $\lambda_i>0$
for all $i$, equivalently $x^\top Sx>0$ for all $x\ne0$.

\cue{Projector} $E_1=\operatorname{range}P$, $E_0=\Null(P)$
for $P^2=P$ (when those eigenvalues occur).
$\tr P=\rank P=\dim\operatorname{range}P$.

\par\vspace{4pt}
{\setlength{\fboxsep}{4pt}%
\noindent\colorbox{wash}{\parbox{\dimexpr\linewidth-2\fboxsep\relax}{%
\raggedright\fontsize{9.2}{11}\selectfont
\setlength{\parskip}{1.6pt}%
{\color{accent}\bfseries Complex numbers and DFT quick tips}\par
\cue{Complex numbers} For $z=a+jb$, $j^2=-1$,
$\bar z=a-jb$, and $\abs z=\sqrt{a^2+b^2}=r$.
Write $z=re^{j\theta}=r(\cos\theta+j\sin\theta)$ and check the quadrant of $\theta$.
For $z^p$, use magnitude $r^p$ and angle $p\theta$.
The $p$ roots use magnitude $\sqrt[p]{r}$ and angles $\frac{\theta+2\pi m}{p}$,
$m=0,\dots,p-1$.
Nonreal eigenvalues of a real matrix come in conjugate pairs.

\cue{DFT / IDFT} In \cue{DFTEXAM}, enter samples in order from $x[0]$.
For $0\le k,n<N$, use $X=Wx$, where $W_{kn}=e^{-jkn\alpha}$ and $\alpha=\frac{2\pi}{N}$.
The inverse flips the exponent sign and scales by $\frac{1}{N}$.
$X[0]=\sum_n x[n]$, so the mean is $\frac{X[0]}{N}$.
For real samples, $X[N-k]=\overline{X[k]}$ for $1\le k<N$.

\cue{Row direction} Wrap the DFT row step $-\frac{2\pi k}{N}$ into $(-\pi,\pi]$.
Negative means clockwise, positive means anticlockwise,
$0$ stays fixed, and $\pi$ alternates.
}}}
\end{panel}

\checkstrip{%
An orthogonal set containing $0$ is dependent.\quad
Repeated eigenvalues can still yield an eigenbasis (as in $I_n$).\quad
A common spectrum alone leaves similarity undecided.}
\end{document}
